% Automatisch aus dem Entwurfs-Doc "Quellsysteme Dart-Händler" (Kap. 6) erzeugt. Anpassen erlaubt.
\begin{landscape}
\begin{small}
\begin{longtable}{@{}r p{3.4cm} p{2.6cm} p{5.2cm} p{7.4cm}@{}}
\caption{Checkliste: Integrationsprobleme aus dem Skript, ihr Einbau in die Quellsysteme und der Power-Query-Schritt}\label{tab:checkliste}\\
\toprule
\textbf{Nr.} & \textbf{Problem (Skript)} & \textbf{Abschnitt} & \textbf{Wo eingebaut} & \textbf{Power-Query-Schritt}\\
\midrule\endfirsthead
\toprule
\textbf{Nr.} & \textbf{Problem (Skript)} & \textbf{Abschnitt} & \textbf{Wo eingebaut} & \textbf{Power-Query-Schritt}\\
\midrule\endhead
\midrule\multicolumn{5}{r}{\itshape Fortsetzung nächste Seite}\\\endfoot
\bottomrule\endlastfoot
1 & Zeichensatz / Codepage & 8.1, 8.2.2.1 synt. & ERP Windows-1252, Shop/CRM UTF-8 & CSV-Import mit Codepage 1252 vs. 65001\\[2pt]
2 & Trennzeichen, Dezimaltrenner & 8.2.2.1 synt. & ERP ; + Komma, Shop , + Punkt & Import-Dialog; Typ ändern mit Gebietsschema de-DE/en-US\\[2pt]
3 & Datumsformate, Zeitzone & 8.2.5 Standardizing & TT.MM.JJJJ / ISO-UTC / MM/DD/YYYY & Typ ändern mit Gebietsschema; DateTimeZone.ToLocal\\[2pt]
4 & Dateiformat & 7.3 Export/Import & txt / csv / xlsx & drei Konnektoren\\[2pt]
5 & Vollabzug vs. Nettoänderung & 7.1 & ERP ohne, Shop mit updated\_at & Filter auf updated\_at \textgreater{} letzter Lauf (diskutieren, nicht zwingend umsetzen)\\[2pt]
6 & Ja/Nein-Kodierungen & 8.2.4 Kodierung & J/N, true/false, X/leer & Bedingte Spalte $\rightarrow$ einheitlich true/false\\[2pt]
7 & Geschlecht/Anrede-Kodierung & 8.2.5, Tab. 8.3 & Shop 1/2/0, CRM Herr/Frau & Werte ersetzen; Widerspruch $\rightarrow$ Priorisierungsregel\\[2pt]
8 & NULL-Kodierungen & 8.1.1.1 & -1, 31.12.9999, leer, 0 & Werte ersetzen $\rightarrow$ null\\[2pt]
9 & Länderkodierung (stark äquivalent) & 8.2.4 & D/A/CH, DE/AT/CH, Klartext & Mapping-Tabelle + Merge\\[2pt]
10 & Kundentyp (schwach äquivalent) & 8.2.4, Abb. 8.17 & 4 Codes / Boolean / 5 Texte & Mapping auf 4 Zielwerte; Großhändler$\rightarrow$Händler (kleinster gemeinsamer Nenner)\\[2pt]
11 & Schematische Heterogenität & 8.2.2.1 & Shop is\_club als Attribut vs. ERP Wert in Spalte & Bedingte Spalte erzeugt Kundentyp aus Boolean\\[2pt]
12 & Attribut vs. Entität & 8.2.2.4, Abb. 8.10 & Lieferart: ERP Spalte, Shop Tabelle & Merge shipping\_methods $\rightarrow$ Spalte; Zielschema: Attribut (begründen!)\\[2pt]
13 & Komposite Schlüssel & Abb. 8.11 & AuftragNr+Pos, order\_id+line\_no & Schlüsselspalte verketten\\[2pt]
14 & Sprechende Schlüssel & 8.2 & ArtNr FL-STD-0421 & Spalte teilen nach Trennzeichen $\rightarrow$ Produktgruppe\\[2pt]
15 & Synonyme & 8.2.2.3, Abb. 8.8 & Vertreter/Betreuer, KdNr/customer\_id/AccountId, Firma/company & Spalten umbenennen (Schemaangleichung)\\[2pt]
16 & Homonyme & 8.2.2.3 & Status $\times$ 3, Segment vs. Kundengruppe & Umbenennen zu Auftragsstatus / Zahlungsstatus / Kontaktstatus\\[2pt]
17 & Fehlende Attribute & 8.2.2.3 & Rolle nur im CRM, E-Mail nicht im ERP & Left Outer Merge, null zulassen und dokumentieren\\[2pt]
18 & Nicht explizit ausgewiesene Attribute & 8.2.2.3 & Shop ohne Verkäufer $\rightarrow$ Kanal = Online & Benutzerdefinierte Spalte Kanal = \grqq{}Online\grqq{} vor dem Append\\[2pt]
19 & Datentypkonflikte & 8.2.2.3, Abb. 8.9 & KdNr Zahl/Text/GUID, PLZ Zahl vs. Text & Alle Schlüssel $\rightarrow$ Text; PLZ $\rightarrow$ Text mit führender Null\\[2pt]
20 & Konfligierende Defaults & 8.2.2.3 & Zahlungsziel Kunde 30 vs. Auftrag 14 & Regel: Auftragswert gewinnt (begründen)\\[2pt]
21 & Referenzielle Integrität & 8.1.1.1 & Aufträge von ausgeschiedenem V01; falsche ERP\_KdNr im CRM & Merge mit Anti-Join $\rightarrow$ Restliste; V01 bleibt mit Austrittsdatum\\[2pt]
22 & Schlüsselbehandlung / Surrogat & 8.2.2.3 & drei Schlüsseltypen & Index-Spalte als KundeKey + Mapping-Tabelle Quell-ID $\leftrightarrow$ Key\\[2pt]
23 & Titel im Namen & Abb. 8.9 & ERP Dr. Müller & Spalte teilen / Text.BeforeDelimiter\\[2pt]
24 & Separierung von Attributen & 8.2.4, 8.2.5 Elementizing & Shop address ein Feld; ERP Straße+Hausnr & Spalte teilen nach Komma / nach letztem Leerzeichen\\[2pt]
25 & Zeichenketten vereinheitlichen & 8.2.4, 8.2.5 Standardizing & Str./Straße/-Straße, ue/ü, e.V. & Werte ersetzen (Liste), Text.Lower, Trim\\[2pt]
26 & Maßeinheiten / Währung & 8.2.4 & CHF$\rightarrow$EUR, brutto$\rightarrow$netto je Land, Set$\rightarrow$Stück & Benutzerdefinierte Spalte mit Kurs-/MwSt-Tabelle\\[2pt]
27 & Record Linkage & 8.2.2.3, 8.2.3 & 130 überlappende Kunden, 3 Schreibweisen & Normalisierter Matchschlüssel (PLZ+Straße+Hausnr) + Fuzzy Merge mit Schwellwert + Transformationstabelle\\[2pt]
28 & Widersprüchliche Werte / Priorisierung & 8.2.3, Tab. 8.3 & Adresse ERP vs. CRM, Anrede & Regel: jüngeres Änderungsdatum gewinnt, sonst ERP (begründen)\\[2pt]
29 & Dubletten innerhalb einer Quelle & 8.1.1.1 Redundanz & CRM 5 \% & Gruppieren nach Matchschlüssel, Anzahl \textgreater{} 1 $\rightarrow$ Prüfliste\\[2pt]
30 & Plausibilität (Verification) & 8.2.5 & PLZ passt nicht zum Land & Bedingte Spalte: Land=CH und Länge(PLZ)$\neq$4 $\rightarrow$ Flag\\[2pt]
31 & Matching / Householding & 8.2.5 & Verein mit mehreren Ansprechpartnern & Gruppieren nach Firma+PLZ\\[2pt]
32 & Filterung & 8.3 & Gesperrt=J, status cancelled/refunded, Leads & Zeilen filtern, Anzahl protokollieren\\[2pt]
33 & Abgeleitete Kennzahlen & 8.3 & Umsatz, Lieferdauer, Zahlungsverzug & Benutzerdefinierte Spalten\\[2pt]
34 & Anreicherung & 8.2.5 & Region aus PLZ, Kanal, Saison aus Datum & Mapping-Tabelle PLZ-Präfix$\rightarrow$Region; Date.QuarterOfYear\\[2pt]
35 & Integrationsstrategie & 8.2.2.2, Abb. 8.7 & leiterartig: ERP $\rightarrow$ +CRM $\rightarrow$ +Shop & Reihenfolge der Merges; Alternative n-är diskutieren\\[2pt]
36 & Mängelklassen & Tab. 8.1 & automatisch / halbautomatisch / manuell & Jede Zeile oben einer Klasse zuordnen (= harte/weiche Regeln aus PA1)\\[2pt]
37 & Historisierung bei Adressänderung & 8.5, Abb. 8.20 & 20 \% Umzieher: ERP alte, CRM neue Adresse mit Änderungsdatum & Beide Versionen behalten: GueltigVon/GueltigBis + CurrentFlag; Auftrag per Datumsbereich zuordnen (Merge + Filter Datum zwischen Von und Bis)\\[2pt]
38 & Zeilenumbrüche und Trennzeichen in Freitext & 8.1 synt. Mängel, 8.2.2.1 & ERP KUNDE.Bemerkung & Csv.Document mit QuoteStyle.Csv statt QuoteStyle.None; danach Text.Replace(\#(cr)\#(lf) $\rightarrow$ \grqq{} | \grqq{}); Vorher/Nachher-Zeilenzahl als Nachweis\\[2pt]
39 & Whitespace: nachgestellt, doppelt, geschützt (NBSP), Nur-Whitespace-Felder & 8.1 synt. Mängel, 8.2.5 Standardizing, Tab. 8.1 Klasse 1 & ERP Firma (trailing), Bemerkung (nur Leerzeichen); Shop address (Doppel-Space); CRM Firma (NBSP, trailing) & fnClean global auf alle Textspalten VOR dem Matchschlüssel; Nachweis: Zahl veränderter Zeilen, Matchtreffer mit/ohne fnClean\\[2pt]
40 & E-Mail-Transportmüll: mailto:, \glqq{}Name \textless{}mail\textgreater{}\grqq{}, GROSS, ;/. am Ende, (at) & 8.1.1.1 fehlerhafte Werte, 8.2.5 Standardizing + Verification & Shop email, CRM E-Mail: je \textasciitilde{}15 \% & fnNormEmail (retten), dann Email\_ok (prüfen, da kein Regex in M – Text.Select/Text.Split); zwei Zahlen: gerettet / weiterhin ungültig\\[2pt]
41 & Fehlbelegung: Hausnummer vor Straße bzw. Felder vertauscht & 8.1.1.1 fehlerhafte Werte, 8.2.5 Verification & ERP Strasse \glqq{}95 Goethestr.\grqq{} (2 Zeilen); CRM Strasse/Hausnr vertauscht (2 Accounts) & Nach dem Teilen: beginnt Strasse mit Ziffer und Hausnr nicht $\rightarrow$ tauschen (bedingte Spalten); Rest wie \glqq{}5 Straße des 17. Juni\grqq{} bleibt manuell (Tab. 8.1 Klasse 2)\\[2pt]
42 & Ungültige Werte: PLZ zu kurz/lang/Buchstabe, E-Mail ohne @ & 8.1.1.1 fehlerhafte Werte & ERP PLZ 2 \%, Shop email 2 \% & Flag\_PLZ\_ok / Email\_ok (Text.Select, Länge) $\rightarrow$ Prüfliste; nicht automatisch korrigierbar (Tab. 8.1 Klasse 2)\\[2pt]
43 & Fehlende Werte: Land leer, PLZ leer & 8.1.1.1 fehlende Werte, 8.1.3 & CRM 3 \% / 2 \% & Land aus PLZ-Länge/Präfix ableiten (wahrscheinlichster Wert, 8.1.3); PLZ leer $\rightarrow$ kein Matchkey, nur Schlüssel/E-Mail/Fuzzy\\[2pt]
44 & Genauigkeit: Datum nur MM/YYYY & 8.1.1.1 Genauigkeit, 8.2.5 Standardtag & CRM Erstkontakt 5 \% & Länge 7 $\rightarrow$ \grqq{}01/\grqq{} voranstellen (Standardtag = 1., Skript 8.2.5) + Flag\_Datum\_ungenau\\[2pt]
45 & Mojibake: ü statt ü (UTF-8 als 1252 gelesen) & 8.2.2.1 synt. Heterogenität Zeichensatz – PA1-Thema & CRM Nachname/Firma 5 \% & Werte ersetzen ü$\rightarrow$ü, ö$\rightarrow$ö, ä$\rightarrow$ä, ß$\rightarrow$ß (Mapping-Liste in fnClean oder eigener Schritt)\\[2pt]
46 & Telefonnummern in 5 Schreibweisen & 8.2.4 Zeichenketten, 8.2.5 Standardizing & ERP Telefon, Shop phone, CRM Telefon & fnTelE164: Ziffern behalten, 00$\rightarrow$+, führende 0 $\rightarrow$ Ländervorwahl aus LandISO (wie PA1); setzt gültiges Land voraus $\rightarrow$ Reihenfolge!\\[2pt]
47 & Dubletten im ERP & 8.1.1.1 Redundanz, 8.2.5 Matching & 3 Paare, andere KdNr, Schreibweise, jüngeres Anlagedatum & Gruppieren nach Matchkey+Name\_norm in KUNDE\_h; ältere KdNr führend, beide auf einen KundeKey (MAP\_SCHLUESSEL zeigt es)\\[2pt]
48 & Anrede-Varianten Hr./Herrn/HERR/herr & 8.2.4 Kodierung & CRM Anrede 10 \% & Text.Lower + Text.Start 2 Zeichen: \grqq{}he\grqq{}/\grqq{}hr\grqq{}$\rightarrow$Herr, \grqq{}fr\grqq{}$\rightarrow$Frau, sonst null\\[2pt]
49 & Konfligierende Integritätsbedingungen auf Beziehungsebene: Attribut an Entität vs. an Beziehung, 1:n vs. n:m & 8.2.2.3, Beispiel 8.2.2.5 (Preis an Waren vs. an bestellt, Abb. 8.14/8.15) & Verkäufer: ERP KUNDE.VertreterNr (je Kunde) vs. AUFTRAG.VertreterNr (je Auftrag) vs. CRM Betreuer am Account und je Aktivität; Preis: ARTIKEL.EKPreis vs. AUFTRAGSPOS.EPreisNetto vs. Shop nur order\_items; Adresse 1:1 vs. historisiert 1:n & Ziel nimmt die allgemeinere Form: Verkäufer am Auftrag (Stammdaten-Vertreter nur Default für Aufträge ohne Angabe), Preis je Position, Adresse mit Gültigkeit. Begründung: Frage 5 (Kanal) nur so korrekt. ER-Skizze ERP/CRM/Ziel in Kap. 2\\[2pt]
\end{longtable}
\end{small}
\end{landscape}
